The 42.79\% initial flip rate suggested that pretraining could already separate digits by orientation. The follow-up compares full-pixel and projected fine-tuning after projected pretraining to measure any return of sensitivity to stripe orientation. After the main study, we \href{https://github.com/JohnR31N/cs126FinalProj/blob/98aa8e4/IDEC/stage/design.md}{recorded a new prediction (98aa8e4)}. With projected fine-tuning fixed, projecting during pretraining should lower the final flip rate.

At $a=0.8$ we trained both pretraining losses from the same random weights and batches for each seed. We fitted K-means to each resulting embedding and forked it into both fine-tuning arms. All four combinations were trained anew for 100 pretraining epochs and 5,600 updates. Changing pretraining also changes the starting centroids.

\begin{center}\small
\begin{tabular}{llrrr}
\toprule
Pretraining & Fine-tuning & Digit ACC (\%) & Flip rate (\%) & Stripe NMI\\
\midrule
Full pixel & AE + K-means & 56.96 $\pm$ 5.46 & 42.79 $\pm$ 12.20 & 0.1823 $\pm$ 0.0600\\
Full pixel & Full pixel & 58.98 $\pm$ 5.35 & 42.22 $\pm$ 12.73 & 0.1815 $\pm$ 0.0639\\
Full pixel & Projected & 59.96 $\pm$ 5.31 & 39.53 $\pm$ 12.68 & 0.1661 $\pm$ 0.0618\\
\addlinespace
Projected & AE + K-means & 67.43 $\pm$ 4.76 & 1.83 $\pm$ 0.31 & 0.0004 $\pm$ 0.0001\\
Projected & Full pixel & 69.97 $\pm$ 5.46 & 4.73 $\pm$ 0.91 & 0.0007 $\pm$ 0.0003\\
Projected & Projected & 69.61 $\pm$ 5.33 & 2.61 $\pm$ 0.54 & 0.0003 $\pm$ 0.0001\\
\bottomrule
\end{tabular}

\end{center}
\smallcap{\textbf{Table 3.} Strong-stripe follow-up, mean $\pm$ sample SD across five seeds. The AE + K-means rows score the two initializations before fine-tuning.}

\textbf{Result.} Holding projected fine-tuning fixed, projected pretraining reduced flips from 39.53\% to 2.61\%. The paired reduction was $36.92\pm12.54$ percentage points and ranged from 21.41 to 53.22 points across seeds. Digit ACC rose from 59.96\% to 69.61\% ($+9.65\pm8.44$ points). One seed lost 1.88 points of accuracy. Every model assigned test digits to all ten clusters.

The difference was already visible immediately after pretraining (42.79\% versus 1.83\% flips). Even after full-pixel fine-tuning, projected pretraining reduced flips by 37.49 points. The two final estimates of the pretraining effect differed by $0.57\pm2.72$ points. Projecting during fine-tuning also helped stability, by 2.69 points after ordinary pretraining and 2.12 after projected pretraining. The ordinary initializations reproduced the archived assignments exactly. Their fine-tuned predictions varied on rerun, so all these contrasts use the four cells of Table~3.
